





\section{Can Fine-Tuning Close the DENSEWORLD Gap?}
\label{sec:vjepa_adaptation_baselines}

Before introducing FactorJEPA, we ask whether conventional adaptation can recover the predictive structure required by DENSEWORLD without reorganizing the JEPA predictor. We evaluate three complementary strategies: LoRA~\cite{hu2022lora}, DoRA~\cite{liu2024dora}, and Auto-RGN~\cite{lee2023surgical}. All use the same raw clips, optimization steps, optimizer, masking policy, and evaluation protocol; only the parameter-update mechanism changes.

Let $x$ denote a video clip; $M_c$ and $M_t$ the context and target masks; $f_\theta$ the online encoder; $\bar f_{\bar\theta}$ the momentum target encoder; and $g_\phi$ the predictor. All baselines retain the executed V-JEPA objective~\cite{bardes2024vjepa} $\mathcal{L}_{\mathrm{JEPA}}=\mathbb{E}_{x}\big\|g_\phi(f_\theta(M_c\odot x),M_t)-\operatorname{sg}\,\bar f_{\bar\theta}(M_t\odot x)\big\|_2^2$, where $\operatorname{sg}$ is stop-gradient. This data- and step-matched protocol isolates adaptation capacity without introducing a different prediction target or training signal.

The gap is real. On an 11-way motion-class probe (19.5\% majority baseline), ten off-the-shelf \emph{frozen} encoders (DINOv2, I-JEPA, LeJEPA, and V-JEPA~1/2.0/2.1) collapse into a narrow 37.5\% to 44.4\% band (Table~\ref{tab:frozen_scorecard}), whereas DENSEWORLD-adapted encoders reach 50.3\% to 53.2\%. Conventional fine-tuning (LoRA, DoRA, Auto-RGN) narrows but does not close the remaining predictive gap, motivating an adaptation that reorganizes \emph{which} factors the predictor attends to rather than only \emph{how many} parameters move.

% Table "Protocol-matched V-JEPA adaptation baselines" (tab:vjepa_adaptation_baselines) relocated to Appendix C to meet the 7-page limit.


Formal baseline update rules, protocol-matched configurations, and resource accounting are in Appendix~C.

\begin{table*}[t]
\centering
\caption{\textbf{\emph{Frozen encoders collapse.}} Ten frozen encoders on the DENSEWORLD motion probe ($n_{\mathrm{test}}=1{,}825$, $\pm$95\% BCa CI): Action top-1 (A) stays in a $37.5$--$44.4\%$ band (19.5\% majority), while adapted encoders reach $50.3$--$53.2\%$. M~=~motion-cos, T~=~taxonomy F1 ($\uparrow$ better); shade $\propto$ column rank.}
\label{tab:frozen_scorecard}
\footnotesize\bfseries
\setlength{\tabcolsep}{4pt}
\renewcommand{\arraystretch}{1.0}
\begin{tabular}{@{}lccc@{\hskip 2.2em}lccc@{}}
\toprule
Frozen encoder & A\,(\%) & M & T & Frozen encoder & A\,(\%) & M & T \\
\midrule
V-JEPA 2.1 (2B) & \cellcolor{green!68}44.4 & \cellcolor{green!36}0.009 & \cellcolor{green!61}0.793 & I-JEPA ViT-H & \cellcolor{green!31}39.1 & \cellcolor{green!58}0.016 & \cellcolor{green!57}0.781 \\
V-JEPA 2.1 ViT-L & \cellcolor{green!67}44.2 & \cellcolor{red!38}0.004 & \cellcolor{green!59}0.788 & V-JEPA 2.0 (SSv2) & \cellcolor{green!30}38.8 & \cellcolor{green!30}0.007 & \cellcolor{green!56}0.776 \\
V-JEPA 1 ViT-H & \cellcolor{green!41}40.5 & \cellcolor{green!30}0.007 & \cellcolor{green!33}0.702 & DINOv2 & \cellcolor{green!27}38.5 & \cellcolor{green!58}0.016 & \cellcolor{green!68}0.816 \\
LeJEPA ViT-L & \cellcolor{green!38}40.1 & \cellcolor{green!52}0.014 & \cellcolor{green!45}0.740 & V-JEPA 2 ViT-L & \cellcolor{green!23}37.9 & \cellcolor{green!49}0.013 & \cellcolor{green!56}0.778 \\
V-JEPA 1 ViT-L & \cellcolor{green!37}39.9 & \cellcolor{green!33}0.008 & \cellcolor{red!38}0.660 & I-JEPA ViT-G/16 & \cellcolor{red!38}37.5 & \cellcolor{green!68}0.019 & \cellcolor{green!59}0.787 \\
\bottomrule
\end{tabular}
\end{table*}




















































% Factor-view figure (factor_visual + segmented) relocated to the end of 2_data.tex (§2 motivation) so it lands on p3.

% Fig 5 (DENSEWORLD vs West density) MOVED to Appendix (app:dataset) to fit the 7-page limit and free room for the frozen-encoder ablation (Table~\ref{tab:frozen_scorecard}). Both the bold bar-plot version and the merged dual-axis line version live there.



\section{\emph{FactorJEPA}: Predictor Surgery over Factorized Views}
\label{sec:factorjepa}

\paragraph{\textbf{\emph{Motivation.}}}
With the same notation, a conventional JEPA predictor $g_\phi$ learns
\begin{equation*}
g_\phi:(f_\theta(M_c\odot x),M_t)\longmapsto\widehat Y_{t+\Delta},
\qquad
Y_{t+\Delta}^{\star}=\operatorname{sg}\!\left[\Pi_{M_t}\bar f_{\bar\theta}(x)\right],
\end{equation*}
by matching $\widehat Y_{t+\Delta}$ to $Y_{t+\Delta}^{\star}$. The objective specifies
\emph{what} to predict, but every clip presents layout, agents and interactions
simultaneously, so nothing constrains \emph{which} of them the predictor actually uses.

This ambiguity is consequential in DENSEWORLD, where layout, agent density,
visibility and interaction pressure are strongly correlated. A high-capacity predictor can
exploit \emph{shortcut mixtures} (crowd texture as a proxy for interaction,
visible appearance for object persistence, or road geometry for motion) without recovering
the structure governing scene evolution.

FactorJEPA addresses this on the \emph{input} side rather than by redesigning
the predictor. We construct three complementary views of every clip, in which layout, agents
and interactions are isolated by construction, and adapt a pretrained V-JEPA predictor over
them as a staged curriculum. The architecture is unchanged; what changes is the sequence of
distributions the predictor is asked to model, and how much of the encoder is free to move at
each stage.

\begin{figure*}[ht!]
    \centering
    \includegraphics[width=\textwidth]
    {figures/forest_plot_best_ci_paper.pdf}
    \caption{\textbf{\emph{FactorJEPA vs. the best competitor, across scale.}} Per-metric advantage in 95\%-CI units at the stratified 10k scale (2B and 1B) and the full 115k scale (1B); green past the dashed $1\times$ line = a separated win. Predictive diagnostics separate and \emph{strengthen} at full scale, where the motion-cosine trade-off also reverses to a win. Full-scale 2B was not run (compute budget).}
    \label{fig:frozen_forest}
    \vspace{-0.6em}
\end{figure*}

\begin{figure*}[t!]
    \centering
    \includegraphics[width=\textwidth]
    {figures/eval_scorecard_winbars.pdf}
    \caption{\textbf{\emph{FactorJEPA outperforms every fine-tuning rival on the four predictive metrics it wins outright}} (Future-frame L1 at 2B; Causal future-block L1, Mask-ratio slope, exposure-bias gap at 1B). Bars show held-out score $\pm$ 95\% BCa CI; every rival is named in the key. Full 15-diagnostic scorecard in the supplementary material. {\footnotesize \emph{wSE}~=~WiSE-FT (weight-space ensembles for fine-tuning).}}
    \label{fig:eval_scorecard}
    \vspace{-0.6em}
\end{figure*}


\subsection{\textbf{\emph{Automatic Factor-View Construction}}}
\label{sec:factor_target_construction}

For each privacy-filtered clip $x_n$ we run an open-vocabulary detector
(Grounding DINO, Swin-B)~\cite{liu2024groundingdino} to localise dynamic agents, and propagate
those detections through the clip with a promptable video segmenter
(Segment Anything Model, SAM~3)~\cite{carion2025sam3}. Detections are re-seeded at four anchor frames and propagated
within each anchor's segment, which keeps identities stable under the heavy occlusion typical
of the regime. This yields, per clip, a set of agent masks and temporally associated tracklets
(a mask and box sequence per agent over its lifespan; formal notation in Appendix~C).
Figure~\ref{fig:dino_factor_targets} illustrates this factor-target construction.

A deterministic map converts these into three \emph{views} of the same clip, each a video the
predictor can be trained on directly (Figure~\ref{fig:factorjepa_visual}):

\begin{itemize}[leftmargin=1.2em,itemsep=2pt,topsep=2pt]
\item $D_L$ (layout): agent regions are suppressed, leaving slowly varying spatial
      support: road geometry, buildings, drivable area.
\item $D_A$ (agents): background context is suppressed, leaving the dynamic actors and
      their motion.
\item $D_I$ (interactions): spatio-temporal tubes around co-occurring, spatially
      proximate tracklet pairs, isolating local multi-agent coupling.
\end{itemize}

Checkpoints, prompt vocabulary, association rules, thresholds and per-view
coverage are reported in Appendix~C.

\subsection{\textbf{\emph{Three-Stage Curriculum}}}

Training proceeds over the three views in order of increasing relational complexity,
$D_L\rightarrow D_A\rightarrow D_I$, with a mixture weight over views at each stage. The
ordering is deliberate: layout is the most stationary signal and anchors the encoder before
agent appearance and finally interaction structure are introduced.

Encoder capacity is released progressively: at stage $s$ only the predictor plus the
deepest $K_s=\lfloor \kappa_s L\rfloor$ transformer blocks (of depth $L$) receive
gradients, so that at most $8/48$ blocks (V-JEPA~2.1 2B) and $6/40$ blocks (V-JEPA~2.1 1B) ever move. The lower blocks and the momentum target network follow the frozen / momentum
protocol of the underlying V-JEPA implementation.

\subsection{\textbf{\emph{Training Objective}}}

FactorJEPA retains the executed V-JEPA objective and adds three auxiliary terms:
\begin{equation*}
\boxed{
\mathcal L=\mathcal L_{\mathrm{JEPA}}
+\beta\,\mathcal L_{\mathrm{probe}}
+\gamma\,\mathcal L_{\mathrm{motion}}
+\lambda\,\mathcal L_{\mathrm{drift}} }
\end{equation*}
where $\mathcal L_{\mathrm{JEPA}}$ is the masked future-latent loss above,
$\mathcal L_{\mathrm{probe}}$ is a lightweight multi-task read-off used for checkpoint
selection, $\mathcal L_{\mathrm{motion}}$ jointly regresses a $13$-dimensional optical-flow
descriptor and classifies its motion class, and $\mathcal L_{\mathrm{drift}}$ is an
$\ell_2$ anchor to the initialisation $\theta^{(0)}$ that bounds how far the adapted encoder
may travel from the pretrained solution.

\subsection{\textbf{\emph{Predictor Surgery}}}

FactorJEPA is initialised from a pretrained V-JEPA model; the online encoder, momentum target
encoder, target construction and masking policy are all retained. Because only the top-$K$
blocks and the predictor are updated, and because $\mathcal L_{\mathrm{drift}}$ penalises
displacement from $\theta^{(0)}$, adaptation is deliberately \emph{local}: it changes how the predictor
is trained, not what it is made of. We refer to this as
\emph{predictor surgery}. Exact stage boundaries, mixture
weights, unfreeze fractions, optimiser settings and measured GPU-hours are reported in
Appendix~C.

\begin{figure}[t!]
    \centering
    \includegraphics[width=0.96\columnwidth]
    {figures/scale_replication_single.pdf}
    \caption{\textbf{\emph{Method rankings survive the 2B$\to$1B backbone change.}} Per-method \emph{Causal future-block L1} at both scales (Spearman $\rho=0.978$), making the 1B backbone a low-cost proxy for 2B. Full 15-panel grid in the supplementary material.}
    \label{fig:scale_replication}
    \vspace{-0.6em}
\end{figure}

\section{Experiments \& Evaluation}
\label{sec:experiments}

We evaluate whether FactorJEPA improves the predictive structure of V-JEPA
rather than only its downstream semantics, using the four diagnostics defined below.


\subsection{\textbf{\emph{Evaluator Independence and Held-out Split}}}
\label{sec:evaluator_independence}

To prevent agreement with the pseudo-label generator from being mistaken for
improved world modeling, factor targets (Grounding DINO + SAM~3) are constructed
only on the training split, and evaluation uses a disjoint held-out split never seen
in training, selection, or tuning. No headline metric reuses FactorJEPA's
pseudo-label generator as its evaluator; each is scored by a frozen, independent
evaluator (full provenance in Appendix~E).

\subsection{\textbf{\emph{Protocol and Comparisons}}}

We evaluate V-JEPA~2.1~\cite{murlabadia2026vjepa21, assran2025vjepa2} at two scales: V-JEPA~2.1 (2B) and V-JEPA~2.1 (1B),
approximately 2B and 1B parameters respectively. The comparison spans the frozen
backbone, conventional full fine-tuning, LP-FT~\cite{kumar2022lpft}, parameter-efficient LoRA, DoRA,
Auto-RGN, FactorJEPA-RAW, and full FactorJEPA.
FactorJEPA-RAW applies the factorized predictor to raw clips without
factor-target supervision, so the chain generic adaptation
(LoRA, DoRA, Auto-RGN) $\to$ FactorJEPA-RAW $\to$ FactorJEPA separates gains due to
parameter adaptation, predictor structure, and structured supervision. All methods
share the same clips, masking policy, evaluation split, and metric implementations;
we report per-encoder test results with 95\% BCa bootstrap confidence intervals.


% Table "Training and evaluation provenance" (tab:evaluator_provenance) relocated to Appendix E to meet the 7-page limit.




\subsection{\textbf{\emph{Four Diagnostics of Predictive Structure}}}
We report four predictive-structure diagnostics (formal definitions in
Appendix~E): \emph{Future-frame L1} ($\downarrow$), the normalized distance
between predicted and target future embeddings; \emph{Causal future-block L1}
($\downarrow$), a strictly causal test that forecasts the target encoder's
future-half latents from the past half alone; \emph{Mask-ratio slope}
($\downarrow$), the growth in prediction error as evidence is progressively
masked; and \emph{Motion cosine} ($\uparrow$), intra- minus inter-class cosine
similarity across motion classes.



\subsection{\textbf{\emph{Decisive Gains over the Frozen World Model}}}

Figure~\ref{fig:frozen_forest} reports FactorJEPA's advantage over the frozen
backbone in confidence-interval widths (values beyond $1\times$ indicate
separation): it clears the threshold on all four headline diagnostics at both
scales, and stays decisive from 2B to 1B. At 1B, FactorJEPA cuts future-frame L1 by 10.6\%, causal future-block L1 by 8.9\%, and mask-ratio slope by 27.4\% against the frozen backbone, and lifts motion-cosine separation to 15.3$\times$ its frozen value, each gain exceeding 70$\times$ the paired 95\% CI. Against the strongest \emph{competing}
adaptation (Figure~\ref{fig:eval_scorecard}) it further separates from \emph{every}
fine-tuning rival on Future-frame L1 and Causal future-block L1 (both scales),
Mask-ratio slope, and free-running exposure-bias gap (1B), for 24 of 24 paired wins over the six baselines, and it beats the strongest full-scale rival, LoRA, on all four headline diagnostics (+2.8\%, +1.6\%, +25.2\%, and +17.8\%). The gain is structured,
not universal: FactorJEPA leads the predictive and masking diagnostics while
frozen and conventionally tuned models keep motion-linearity and some
frame-timing~\cite{misra2016shuffle, wang2020pace} and
temporal-coherence~\cite{wei2018aot, dwibedi2019tcc} measures. The lone
trade-off is Motion cosine, where factorization favours accurate, causally
structured prediction over maximally linear motion readout (per-metric
magnitudes in Appendix~D).

\subsection{\textbf{\emph{Cross-Scale Replication}}}
The complete ranking of adaptation methods is stable across the 2B and 1B backbones (20 arms at 2B, 14 at 1B).
The per-diagnostic Spearman correlations $\rho_k$ are strongest on the four headline
metrics ($\rho=0.895$ to $0.978$; Figure~\ref{fig:scale_replication}), twelve of
fifteen diagnostics retain consistent rankings, and only Temporal-order, Teacher-free
gap, and Rollout drift fail to transfer (full comparison in
the supplementary material). This establishes the 1B model as a faithful
lower-cost proxy for 2B, and shows the gains and the motion trade-off are
properties of the method ordering, not artifacts of one backbone.


\subsection{\textbf{\emph{Ablations}}}

We isolate each ingredient with matched controls on the same held-out split (paired
BCa bootstrap; \emph{separates} means the oriented 95\% interval excludes zero); full
per-metric detail is in Appendix~D.
\textbf{\emph{The gain is predictor surgery, not the factor targets:}}
FactorJEPA-RAW (identical schedule and unfreeze budget, no factor supervision) is not
beaten by full FactorJEPA on future-frame at either scale, so the closing of the
frozen-to-adapted gap comes from the surgical adaptation, not the factor-view
supervision layered on top.
\textbf{\emph{Updating the encoder, not only the head, drives the wins:}}
head-only adaptation loses on future-frame, causal future-block, and mask-ratio (each
separated) but wins back motion-cosine, rollout, and exposure-bias.
\textbf{\emph{The interaction stage adds no measurable value:}} removing $D_I$
(no-$D_I$) leaves every predictive diagnostic statistically unchanged.
\textbf{\emph{Weight merging hurts:}} a WiSE-FT (weight-space ensembles for fine-tuning)~\cite{wortsman2022wiseft} interpolation (mixing fractions
$0.3$, $0.5$, $0.7$) loses to the un-merged model on future-frame and motion-cosine at
both scales.

\paragraph{\textbf{\emph{Limitations.}}}
Three controls bound our claims. A compute-matched control (plain pretraining for the
same number of steps as the surgery) is declared but not run, so we cannot fully
separate the surgical schedule from additional optimisation; head-only pretraining and
the continual-learning baselines (CaSSLe~\cite{fini2022cassle}, EWC (elastic weight consolidation)~\cite{kirkpatrick2017ewc}) were trained but not evaluated on the
suite; and the method is a precondition on the base checkpoint, as on V-JEPA~2.0 the
same recipe does not improve over the frozen encoder, so the gains are specific to
V-JEPA~2.1. At the 10k screening scale factorization also concedes maximally linear
motion readout to full fine-tuning ($-44.9$\% at 2B, $-33.4$\% at 1B), a trade-off
that reverses only at full ($86{,}831$-video) scale.













































